\section{Local variation of the complete physical remainder}
\label{sec:v49-physical-variation}
The multiplier estimate applies directly to the source left in
Section~\ref{sec:v48-physical-remainder}. We sum the actual physical
margins before taking a collision limit. The resulting estimate is
local in the roof variable and retains every exact discrete label.

For $f\in L^1_{\rm loc}(\R)$ and $h>0$ define
\begin{equation}\label{eq:v49-local-variation-norm}
 \|f\|_{1,h}^{\rm loc}
       =\sup_{t\in\R}\int_t^{t+h}|f(u)|\,du.
\end{equation}
For a signed or complex measure this notation refers to its total
variation on the interval. No factor $1/2$ occurs in this norm.
The factor $1/2$ will be used only for the distance between two
probability laws.

Keep the physical guard $H_{m,R}^\varepsilon$ of
\eqref{eq:v48-physical-remainder}, and put
$q_0=47/53$, $x_0=q_0^{1/64}$. The two depth functions are
$d_j=\min(j,m-j)$ for contacts and
$d'_j=\min(j,m-1-j)$ for flights. A physical factor below one has,
respectively,
\[
 c_j<2\varepsilon q_0^{d_j/4}
 \quad\hbox{or}\quad
 \Delta_j<2\varepsilon q_0^{d'_j/4}.
\]
We initially take $0<2\varepsilon<s_0$; increasing constants handles
any fixed remaining compact range of margins.

\begin{theorem}[A local bound on all physical defects]
\label{thm:v49-physical-remainder-local}
Let $|w|\le M$, with no regularity assumption on $w$. For every fixed
auxiliary band $B_*\ge B_0$, every $h>0$ and $m\ge2$,
\begin{align}
 &\sup_{R,n,k}m^2
        \|b^{\varepsilon,w}_{n,k,m,R}\|_{1,h}^{\rm loc}\notag\\
 &\qquad\le CM\varepsilon^{1/16}(h+B_*^{-1})
       +C_{B_*}M(1+B_*h)m^3\rho_{B_*}^{m/2}.
                                      \label{eq:v49-physical-window-bound}
\end{align}
The same estimate applies to either first-physical-defect type, and
their sum is the original physical remainder. The width $\varepsilon$
may depend on $m$ in this finite-count estimate.

Let $b_{>J}^{\varepsilon,w}$ be the part whose first physical factor
strictly below one has depth greater than $J$, in the physical-only
order used in \eqref{eq:v48-two-physical-types}. Then
\begin{align}
 &\sup_{R,n,k}m^2\|b_{>J,n,k,m,R}^{\varepsilon,w}\|_{1,h}^{\rm loc}
 \le CM\varepsilon^{1/16}x_0^{J+1}(h+B_*^{-1})\notag\\
 &\hspace{37mm}+C_{B_*}M(1+B_*h)m^3\rho_{B_*}^{m/2}.
                                      \label{eq:v49-physical-depth-tail}
\end{align}
An empty tail is zero. In particular, at each fixed margin and roof
length,
\begin{equation}\label{eq:v49-local-limsup}
 \limsup_{m\to\infty}\sup_{R,n,k,|w|\le M}
       \frac{m^2}{h}\|b^{\varepsilon,w}_{n,k,m,R}\|_{1,h}^{\rm loc}
                                    \le CM\varepsilon^{1/16}.
\end{equation}
\end{theorem}
\begin{proof}
For factors in $[0,1]$, $1-\prod_i a_i\le\sum_i\one_{\{a_i<1\}}$.
Thus the original positive source $1-H_{m,R}^\varepsilon$ is bounded
by the union sum of the contact and flight events displayed above.
There are at most two contacts and two flights of any given depth.
Use the incidence indicator at time $j$ and, by
\eqref{eq:v49-clearance-inclusion}, the backward-clearance indicator
at time $j+1$. These marks always lie between zero and $m$.

On the exact return component $\Pi_{n,k,m,R}$, the collision labels
are exactly $(K_m^{\rm c},A_m)=(k,n)$ and the roof is $S_m$.
Enlarging the source from the initial section to the whole cylinder,
and dropping the terminal membership only on the upper side, costs
exactly $1/c$. Apply Theorem~\ref{thm:v49-physical-marked-local} to
every term. The leading terms sum because
\[
 \sum_{a\ge0}(2\varepsilon q_0^{a/4})^{1/16}
                  =\frac{(2\varepsilon)^{1/16}}{1-x_0}.
\]
There are $2m+1$ contact and flight terms in total. Their finite-count
errors sum to $C_{B_*}(1+B_*h)m^3\rho_{B_*}^{m/2}$, before any
limit is taken. This proves \eqref{eq:v49-physical-window-bound}
for the unweighted positive source.

A first physical defect of depth greater than $J$ is contained in
the same union restricted to depths $a>J$. The leading geometric
sum is multiplied by $x_0^{J+1}$; the error can be bounded by the
same full finite sum. This proves \eqref{eq:v49-physical-depth-tail}.
Grouping a disjoint first-defect partition by type gives the two
individual conclusions. We have not identified first-defect pieces
with independently normalized laws.

For complex $w$, pushforward and total variation give
$|b^{\varepsilon,w}(u)|\le M b^{\varepsilon,1}(u)$ almost everywhere.
Absolute continuity is already available for the complete fixed-count
source from Theorem~\ref{thm:v44-complete-density-bound}; no derivative
of $w$ is used. Finally fix $\varepsilon,h,B_*$, take the collision
limsup, and only then let $B_*\to\infty$. The leading constant in
\eqref{eq:v49-physical-window-bound} is independent of $B_*$.
This proves \eqref{eq:v49-local-limsup}. The finite-count assertion
allows $\varepsilon=\varepsilon_m$ because its remainder is independent
of the strip widths; it says nothing about a count-dependent pointwise
protected inverse.
\end{proof}

\subsection{The full signed correction in a local density norm}
The elementary convolution estimate needed here is
\begin{equation}\label{eq:v49-convolution-local}
 \|K*f\|_{1,h}^{\rm loc}\le\|K\|_1\|f\|_{1,h}^{\rm loc}.
\end{equation}
Indeed, integrate the pointwise convolution upper bound over a
translated interval and use Tonelli followed by the supremum in
\eqref{eq:v49-local-variation-norm}. This argument does not require
that $K$ be positive or compactly supported.

\begin{corollary}[All reconstruction bands on the physical source]
\label{cor:v49-physical-all-band}
For every fixed $\varepsilon,h>0$,
\begin{equation}\label{eq:v49-physical-all-band}
 \limsup_{m\to\infty}\sup_{R,n,k,|w|\le M,\,B>0}
 \frac{m^2}{h}\|b^{\varepsilon,w}_{n,k,m,R}
                -K_B*b^{\varepsilon,w}_{n,k,m,R}\|_{1,h}^{\rm loc}
                                 \le CM\varepsilon^{1/16}.
\end{equation}
There is also a version of \eqref{eq:v49-physical-depth-tail} with
the same geometric tail and finite-count error for this correction.
\end{corollary}
\begin{proof}
Apply \eqref{eq:v49-convolution-local} and
$\|K_B\|_1=\|K_1\|_1$ before taking any supremum or limit. Multiply
\eqref{eq:v49-physical-window-bound} or its depth-tail version by
$1+\|K_1\|_1$, take the collision limsup, and then enlarge the
auxiliary band $B_*$. The reconstruction band $B$ and the auxiliary
spectral band $B_*$ are distinct throughout.
\end{proof}

For a protected-admissible insertion $w$, let $(M,L)$ be the supremum
and endpoint-gradient budget of
Theorem~\ref{thm:v46-protected-correction}, uniform in the parameters
being considered. The original insertion $w=1$ has $(M,L)=(1,0)$.
Write $p^w$ for its full exact-label density, with no physical or
section guard left in that density.

\begin{theorem}[The complete raw correction in local variation]
\label{thm:v49-full-local-correction}
Choose $\varepsilon(B)=A B^{-1/12}$ with $A$ sufficiently large as
in Theorem~\ref{thm:v48-two-source-reduction}. For every sufficiently
large fixed $B$ and every fixed $h>0$,
\begin{equation}\label{eq:v49-full-local-correction}
 \limsup_{m\to\infty}\sup_{R,n,k,w}
 \frac{m^2}{h}\|p^w_{n,R}(k,m,\cdot)
             -K_B*p^w_{n,R}(k,m,\cdot)\|_{1,h}^{\rm loc}
                              \le C(M+L)B^{-1/192}.
\end{equation}
The constant does not depend on $h$. The collision threshold and
auxiliary spectral constants may depend on $h$ and the fixed $B$.
Every physical and section source is included on the left.
\end{theorem}
\begin{proof}
Use the exact source identity, also for the same insertion,
\[
 p^w=f^{\varepsilon,w}
           +d^{\varepsilon,\varepsilon,w}+b^{\varepsilon,w}.
\]
The smoothly protected correction has essential-supremum limsup
$C(M+L)B^{-1/2}$ after normalization by $m^2$, since
$B\ge C_2\varepsilon(B)^{-12}$ once $A^{12}\ge C_2$.
Its local variation norm is at most $h$ times that bound.
The all-depth decision correction has normalized essential height
$CM\varepsilon(B)^{1/16}$ by
Corollary~\ref{cor:v48-decision-all-band}; pay it in the same way.
For the physical term apply Corollary~\ref{cor:v49-physical-all-band}.
This last term is estimated in local variation, not in essential
supremum. Now
$\varepsilon(B)^{1/16}=A^{1/16}B^{-1/192}$, and the first term is
smaller. Adding these three estimates proves the theorem.

For each fixed $B$, its margin is fixed before the collision limit.
The auxiliary band $B_*$ in the physical estimate is enlarged only
after that limit, with $B$ still fixed. Only after all these estimates
may $B$ tend to infinity. No spectral estimate has been used at a
band depending on $m$.
\end{proof}

\paragraph{Pointwise values are a separate question.}
Equation~\eqref{eq:v49-full-local-correction} estimates the integral
of the absolute density error, and not only the signed error of a
fixed test interval. It supplies no essential-height bound for the
physical remainder. Even a bounded density of arbitrarily small
local integral can have thin spikes. The physical criterion
\eqref{eq:v48-physical-residual-criterion} therefore retains its
original pointwise meaning and is not declared proved by this theorem.
